\section{Planteamiento del problema}
ThingsBoard se posiciona como una plataforma IoT vers\'atil, pero su desempe\~no real depende de la topolog\'ia de red, de la configuraci\'on de brokers MQTT y de la forma en que evoluciona la carga de telemetr\'ia. La organizaci\'on dispone de una red HaLow que permite aislar el tr\'afico y degradar enlaces de forma controlada, pero carec\'ia de evidencia cuantitativa sobre cu\'antas conexiones simult\'aneas puede sostener el servidor antes de incumplir los indicadores de calidad de servicio (QoS) negociados con los usuarios finales.

El problema se traduce en responder preguntas como: \textit{¿a partir de cu\'antos dispositivos concurrentes la latencia media supera el umbral de \SI{10}{\milli\second}? ¿Qu\'e combinaciones de intervalo de publicaci\'on y rampas generan p\'erdidas o errores MQTT? ¿C\'omo inciden las variaciones de la red HaLow en la estabilidad de ThingsBoard?} Sin un laboratorio reproducible era imposible comparar escenarios o justificar inversiones en hardware/virtualizaci\'on. Este trabajo propone un banco de pruebas que automatiza la generaci\'on de telemetr\'ia y el registro de evidencias para delimitar la regi\'on de operaci\'on segura del servidor.
